% Options for packages loaded elsewhere
\PassOptionsToPackage{unicode}{hyperref}
\PassOptionsToPackage{hyphens}{url}
\documentclass[11pt]{article}
\usepackage{amsmath,amssymb}
\usepackage{iftex}
\ifPDFTeX
  \usepackage[T1]{fontenc}
  \usepackage[utf8]{inputenc}
  \usepackage{textcomp} % provide euro and other symbols
\else % if luatex or xetex
  \usepackage{unicode-math} % this also loads fontspec
  \defaultfontfeatures{Scale=MatchLowercase}
  \defaultfontfeatures[\rmfamily]{Ligatures=TeX,Scale=1}
\fi
\usepackage{lmodern}

% Custom modern typography and page layout
\usepackage[letterpaper, margin=1in]{geometry}
\usepackage{xcolor}
\usepackage{colortbl}
\usepackage{fancyhdr}
\usepackage{titlesec}
\usepackage{enumitem}
\usepackage{pdflscape}

% Color palette definition
\definecolor{DeepNavy}{HTML}{1A365D}
\definecolor{SlateBlue}{HTML}{2B6CB0}
\definecolor{Charcoal}{HTML}{2D3748}
\definecolor{LightGray}{HTML}{F7FAFC}
\definecolor{TableLine}{HTML}{CBD5E0}
\definecolor{AccentGold}{HTML}{D69E2E}

% Modern section styling
\titleformat{\section}
  {\color{DeepNavy}\normalfont\Large\bfseries}
  {\thesection}{1em}{}[{\color{DeepNavy}\titlerule[0.8pt]}]
\titleformat{\subsection}
  {\color{SlateBlue}\normalfont\large\bfseries}
  {\thesubsection}{1em}{}
\titleformat{\subsubsection}
  {\color{Charcoal}\normalfont\normalsize\bfseries}
  {\thesubsubsection}{1em}{}

% Adjust section numbering and TOC depth
\setcounter{secnumdepth}{3}
\setcounter{tocdepth}{3}

% Headers and Footers configuration
\pagestyle{fancy}
\fancyhf{}
\fancyhead[L]{\color{Charcoal}\small\itshape Plan de Negocio}
\fancyhead[R]{\color{Charcoal}\small Eco Logistic S.A.C.}
\fancyfoot[C]{\color{Charcoal}\thepage}
\renewcommand{\headrulewidth}{0.4pt}
\renewcommand{\footrulewidth}{0pt}

% Keep other pandoc package configurations for compatibility
\IfFileExists{upquote.sty}{\usepackage{upquote}}{}
\IfFileExists{microtype.sty}{% use microtype if available
  \usepackage[]{microtype}
  \UseMicrotypeSet[protrusion]{basicmath} % disable protrusion for tt fonts
}{}

\makeatletter
\@ifundefined{KOMAClassName}{% if non-KOMA class
  \IfFileExists{parskip.sty}{%
    \usepackage{parskip}
  }{% else
    \setlength{\parindent}{0pt}
    \setlength{\parskip}{6pt plus 2pt minus 1pt}}
}{% if KOMA class
  \KOMAoptions{parskip=half}}
\makeatother

\usepackage{longtable,booktabs,array}
\newcounter{none} % for unnumbered tables
\usepackage{multirow}
\usepackage{calc} % for calculating minipage widths
\usepackage{etoolbox}
\makeatletter
\patchcmd\longtable{\par}{\if@noskipsec\mbox{}\fi\par}{}{}
\makeatother

\IfFileExists{footnotehyper.sty}{\usepackage{footnotehyper}}{\usepackage{footnote}}
\makesavenoteenv{longtable}
\usepackage{graphicx}

\makeatletter
\newsavebox\pandoc@box
\newcommand*\pandocbounded[1]{% scales image to fit in text height/width
  \sbox\pandoc@box{#1}%
  \Gscale@div\@tempa{\textheight}{\dimexpr\ht\pandoc@box+\dp\pandoc@box\relax}%
  \Gscale@div\@tempb{\linewidth}{\wd\pandoc@box}%
  \ifdim\@tempb\p@<\@tempa\p@\let\@tempa\@tempb\fi% select the smaller of both
  \ifdim\@tempa\p@<\p@\scalebox{\@tempa}{\usebox\pandoc@box}%
  \else\usebox{\pandoc@box}%
  \fi%
}
\def\fps@figure{htbp}
\makeatother

\setlength{\emergencystretch}{3em} % prevent overfull lines
\providecommand{\tightlist}{%
  \setlength{\itemsep}{0pt}\setlength{\parskip}{0pt}}
\usepackage{bookmark}
\IfFileExists{xurl.sty}{\usepackage{xurl}}{} % add URL line breaks if available
\urlstyle{same}
\hypersetup{
  hidelinks,
  pdfcreator={LaTeX via pandoc}}

\author{}
\date{}

\begin{document}

% Portada Corporativa Premium
\begin{titlepage}
  \thispagestyle{empty}
  \centering
  
  \vspace*{0.5cm}
  {\scshape\Large\bfseries UNIVERSIDAD PERUANA UNIÓN \par}
  {\scshape\large Facultad de Ingeniería y Arquitectura \par}
  {\scshape\normalsize Escuela Profesional de Ingeniería de Sistemas \par}
  
  \vspace{0.8cm}
  
\includegraphics[width=3.2in,keepaspectratio]{media/image1.png}
  \vspace{0.8cm}
  
  {\color{DeepNavy}\hrule height 2pt}
  \vspace{0.4cm}
  {\Huge\bfseries PLAN DE NEGOCIO \par}
  \vspace{0.4cm}
  {\color{AccentGold}\hrule height 1pt}
  \vspace{0.4cm}
  {\Large\bfseries ``Eco Logistic S.A.C.'' \par}
  \vspace{0.4cm}
  {\color{DeepNavy}\hrule height 2pt}
  
  \vspace{0.8cm}
  
  \begin{minipage}[t]{0.55\textwidth}
    \begin{flushleft} \large
      \textbf{Integrantes:}\\
      \smallskip
      \small
      Jhack Quispe Mamani\\
      Alvaro Kenji Cuti Magaño\\
      Abel Is-boset Guevara Huasco\\
      Javier Kevin Tello Rivas\\
      David Edison Guerra Ordoñez\\
      Daniel Hernana Calderón Camacho\\
      Richard Josef Yanqui Inocente\\
      Sebastian Chinchay Collachagua\\
      Joshua Jurguen Alfonzo Javier
    \end{flushleft}
  \end{minipage}
  \hfill
  \begin{minipage}[t]{0.35\textwidth}
    \begin{flushright} \large
      \textbf{Curso:}\\
      \smallskip
      \small Gobierno de TI\\
      \vspace{0.4cm}
      \textbf{Docente:}\\
      \smallskip
      \small Mg. Lennin H. Centurión J.
    \end{flushright}
  \end{minipage}
  
  \vfill
  
  {\large Lima --- Perú \par}
  {\large 2025 \par}
\end{titlepage}

\newpage
\pagenumbering{roman}
\thispagestyle{plain}
\tableofcontents
\newpage

\pagenumbering{arabic}
\setcounter{page}{1}

\section{Situación de la empresa}

\subsection{Descripción de la empresa}

Eco Logistic SAC es una empresa peruana con 15 años de trayectoria en el
sector de almacenamiento y logística, bajo la dirección estratégica de
Óscar Neira. La organización se distingue por un modelo de negocio dual
que combina el arrendamiento de espacios industriales con servicios
especializados de operación logística, lo que le permite adaptarse a
diferentes necesidades del mercado y generar valor agregado a sus
clientes.

En cuanto a su modelo de negocio, la empresa desarrolla sus actividades
bajo dos modalidades principales. Por un lado, se enfoca en el
subarrendamiento de espacios, identificando grandes locales industriales
estratégicos que son alquilados en su totalidad para posteriormente ser
subarrendados por partes a distintos clientes mediante contratos de
largo plazo, generalmente entre 3 a 5 años. Por otro lado, Eco Logistic
SAC opera como un operador logístico integral, gestionando directamente
la mercadería de sus clientes. Este servicio incluye procesos como la
recepción, almacenamiento, control de inventarios, picking, embalaje,
reempaque y despacho de productos. Un ejemplo representativo de esta
operación es la gestión de la campaña de panetones para la empresa
Bimbo, donde se movilizan aproximadamente 5,000 paletas mensuales
durante el periodo de junio a diciembre. Respecto a su infraestructura,
la empresa cuenta con una sede de aproximadamente 40,000 metros
cuadrados ubicada en el Callao. Este espacio, que anteriormente
funcionaba como una fábrica de cables, fue incorporado a sus operaciones
tras el periodo de la pandemia. El local dispone de áreas destinadas al
almacenamiento con racks, zonas de despacho y espacios administrativos,
lo que permite una operación logística a gran escala.

En relación con su cartera de clientes, Eco Logistic SAC ha logrado
consolidar alianzas con importantes empresas corporativas,
principalmente del sector de alimentos y bebidas, tales como Bimbo,
Pepsi, Backus, Cola Real (Añaños) y La Calera, además de distribuidoras
de productos reconocidos. Asimismo, ha trabajado con empresas de otros
sectores como Industrias San Miguel (ISM) y compañías de plásticos, e
incluso ha brindado servicios específicos como el almacenamiento de
unidades vehiculares (Crocotaxis), lo que evidencia su capacidad de
adaptación a diferentes tipos de requerimientos logísticos.

No obstante, la empresa enfrenta actualmente un importante desafío en el
ámbito tecnológico. Su sistema anterior, denominado NDS, se encuentra
inactivo y obsoleto, ya que fue desarrollado sin documentación ni
soporte técnico adecuado. Como consecuencia, la organización ha tenido
que retornar temporalmente a procesos manuales, registrando la
información en papel para luego transcribirla a Excel o sistemas
básicos, lo que genera problemas significativos en la trazabilidad y
ubicabilidad de los productos, así como un mayor riesgo de errores
operativos.

En cuanto a su visión estratégica, Eco Logistic SAC tiene como objetivo
principal la profesionalización de sus servicios logísticos. Entre sus
proyectos futuros destaca la intención de convertirse en un terminal de
almacenamiento aduanero, con capacidad para gestionar entre 600 y 700
contenedores bajo la supervisión de la Superintendencia Nacional de
Aduanas y de Administración Tributaria (SUNAT). Asimismo, la empresa
proyecta expandir su modelo de negocio hacia el mercado internacional,
específicamente en Estados Unidos, aprovechando las oportunidades
económicas y la escala de dicho mercado.

Finalmente, es importante señalar que la empresa opera dentro de un
contexto complejo en términos de seguridad, debido a la alta incidencia
de delincuencia en la zona del Callao. Por esta razón, mantiene una
política de bajo perfil, evitando la colocación de letreros externos en
sus instalaciones. Complementariamente, cuenta con medidas de seguridad
como personal armado y sistemas de videovigilancia, con el fin de
proteger tanto la mercadería de sus clientes como la integridad de su
personal.

\subsection{Misión}

``Brindar soluciones logísticas integrales mediante servicios de
almacenamiento, gestión de inventarios y despacho, apoyados en el uso de
tecnología, con el fin de garantizar eficiencia, seguridad y
trazabilidad, generando valor para nuestros clientes.''

\subsection{Visión}

``Ser una empresa logística líder a nivel nacional con proyección
internacional, reconocida por su eficiencia operativa, uso de tecnología
avanzada y capacidad para gestionar terminales de almacenamiento bajo
estándares de calidad y seguridad.''

\subsection{Alineamientos estratégicos}

El plan estratégico de Cologistic S.A.C. para el presente periodo se
enfoca en la consolidación de su infraestructura logística y la
modernización de su gestión operativa. Dado que el modelo de negocio se
basa en el arrendamiento de grandes extensiones y servicios de valor
agregado para corporaciones de alto volumen, los objetivos están
orientados a maximizar la rentabilidad del metro cuadrado y garantizar
una trazabilidad impecable.

Esta planificación estratégica busca la transición de una gestión manual
basada en hojas de cálculo hacia una operación automatizada y segura,
alineando la capacidad física instalada en el Callao con las demandas de
clientes de alta exigencia como Bimbo o ISM

\begin{landscape}
\footnotesize
\setlength{\tabcolsep}{3pt}
\renewcommand{\arraystretch}{1.6}
\arrayrulecolor{TableLine}
{\def\LTcaptype{none} % do not increment counter
\begin{longtable}[]{@{}
  >{\raggedright\arraybackslash}p{(\linewidth - 30\tabcolsep) * \real{0.1170}}
  >{\raggedright\arraybackslash}p{(\linewidth - 30\tabcolsep) * \real{0.1286}}
  >{\raggedright\arraybackslash}p{(\linewidth - 30\tabcolsep) * \real{0.1277}}
  >{\raggedright\arraybackslash}p{(\linewidth - 30\tabcolsep) * \real{0.0374}}
  >{\raggedright\arraybackslash}p{(\linewidth - 30\tabcolsep) * \real{0.0353}}
  >{\raggedright\arraybackslash}p{(\linewidth - 30\tabcolsep) * \real{0.0396}}
  >{\raggedright\arraybackslash}p{(\linewidth - 30\tabcolsep) * \real{0.0342}}
  >{\raggedright\arraybackslash}p{(\linewidth - 30\tabcolsep) * \real{0.0432}}
  >{\raggedright\arraybackslash}p{(\linewidth - 30\tabcolsep) * \real{0.0351}}
  >{\raggedright\arraybackslash}p{(\linewidth - 30\tabcolsep) * \real{0.0342}}
  >{\raggedright\arraybackslash}p{(\linewidth - 30\tabcolsep) * \real{0.0370}}
  >{\raggedright\arraybackslash}p{(\linewidth - 30\tabcolsep) * \real{0.0342}}
  >{\raggedright\arraybackslash}p{(\linewidth - 30\tabcolsep) * \real{0.0342}}
  >{\raggedright\arraybackslash}p{(\linewidth - 30\tabcolsep) * \real{0.0368}}
  >{\raggedright\arraybackslash}p{(\linewidth - 30\tabcolsep) * \real{0.0342}}
  >{\raggedright\arraybackslash}p{(\linewidth - 30\tabcolsep) * \real{0.1912}}@{}}
\toprule\noalign{}
\endhead
\bottomrule\noalign{}
\endlastfoot
\rowcolor{DeepNavy}
\multirow{2}{=}{\textbf{\textcolor{white}{Objetivo Estratégico}}} &
\multirow{2}{=}{\textbf{\textcolor{white}{Indicadores}}} &
\multirow{2}{=}{\textbf{\textcolor{white}{Responsable}}} &
\multicolumn{12}{>{\raggedright\arraybackslash}p{(\linewidth - 30\tabcolsep) * \real{0.4355} + 22\tabcolsep}}{%
\textbf{\textcolor{white}{Cronograma}}} & \multirow{2}{=}{\textbf{\textcolor{white}{Presupuesto}}} \\
\rowcolor{DeepNavy}
& & & \textbf{\textcolor{white}{ene}} & \textbf{\textcolor{white}{feb}} & \textbf{\textcolor{white}{mar}} & \textbf{\textcolor{white}{abr}} &
\textbf{\textcolor{white}{may}} & \textbf{\textcolor{white}{jun}} & \textbf{\textcolor{white}{jul}} & \textbf{\textcolor{white}{ago}} & \textbf{\textcolor{white}{set}}
& \textbf{\textcolor{white}{oct}} & \textbf{\textcolor{white}{nov}} & \textbf{\textcolor{white}{dic}} \\ \hline
Mejorar la rentabilidad de la empresa & Rentabilidad sobre Ventas & Jefe
de Comercialización & X & X & X & X & X & X & X & X & X & X & X & X &
8000 \\ \hline
\rowcolor{LightGray}
Reducción de costos Operativos & Costo Operativo por Unidad Producida
(COUP) & Junta Administrativa & X & & & & & & X & & & & & & 3000 \\ \hline
Creación de nuevos productos para abarcar mayor parte en el mercad &
Participación de nuevos productos en ventas totales & Junta
Administrativa & & & & & X & X & & & & & X & X & 2500 \\ \hline
\rowcolor{LightGray}
Aumento de la visibilidad en linea y captación de nuevos clientes &
Porcentaje de conversión digital & Junta Administrativa & & & & & & X &
& & & & & X & 2000 \\ \hline
Implementar un sistema de monitoreo de calidad. & Porcentaje de
productos con defectos detectados en control de calidad & Ing. de
Calidad (contratar) & X & & & & & & & & & & & & 2000 \\ \hline
\rowcolor{LightGray}
Expansión de flota comercial. para llegar a mas clientes & Índice de
cumplimiento de entregas & Jefe de Comercialización & & X & & & X & & &
X & & & X & & 3000 \\ \hline
Mejorar la gestión de materia prima. & Tasa de desperdicio de insumos &
Jefe de Almacenamiento & X & X & X & X & X & X & X & X & X & X & X & X &
1000 \\ \hline
\rowcolor{LightGray}
Establecer alianzas estratégicas para evitar dependencia de contratos
gubernamentales. & Nivel de concentración de compras & Junta
Administrativa & X & & & & & & & & & & & & 800 \\ \hline
Capacitar al personal & Rentabilidad semanal & Jefe de cada Área & X & &
& X & & & X & & & X & & & 2500 \\ \hline
& & &
\multicolumn{2}{>{\raggedright\arraybackslash}p{(\linewidth - 30\tabcolsep) * \real{0.0727} + 2\tabcolsep}}{%
} &
\multicolumn{2}{>{\raggedright\arraybackslash}p{(\linewidth - 30\tabcolsep) * \real{0.0738} + 2\tabcolsep}}{%
} &
\multicolumn{2}{>{\raggedright\arraybackslash}p{(\linewidth - 30\tabcolsep) * \real{0.0783} + 2\tabcolsep}}{%
} &
\multicolumn{2}{>{\raggedright\arraybackslash}p{(\linewidth - 30\tabcolsep) * \real{0.0712} + 2\tabcolsep}}{%
} &
\multicolumn{2}{>{\raggedright\arraybackslash}p{(\linewidth - 30\tabcolsep) * \real{0.0684} + 2\tabcolsep}}{%
} &
\multicolumn{2}{>{\raggedright\arraybackslash}p{(\linewidth - 30\tabcolsep) * \real{0.0710} + 2\tabcolsep}}{%
} & \textbf{24800} \\ \hline
& & &
\multicolumn{2}{>{\raggedright\arraybackslash}p{(\linewidth - 30\tabcolsep) * \real{0.0727} + 2\tabcolsep}}{%
} &
\multicolumn{2}{>{\raggedright\arraybackslash}p{(\linewidth - 30\tabcolsep) * \real{0.0738} + 2\tabcolsep}}{%
} &
\multicolumn{2}{>{\raggedright\arraybackslash}p{(\linewidth - 30\tabcolsep) * \real{0.0783} + 2\tabcolsep}}{%
} &
\multicolumn{2}{>{\raggedright\arraybackslash}p{(\linewidth - 30\tabcolsep) * \real{0.0712} + 2\tabcolsep}}{%
} &
\multicolumn{2}{>{\raggedright\arraybackslash}p{(\linewidth - 30\tabcolsep) * \real{0.0684} + 2\tabcolsep}}{%
} &
\multicolumn{2}{>{\raggedright\arraybackslash}p{(\linewidth - 30\tabcolsep) * \real{0.0710} + 2\tabcolsep}}{%
} & \textbf{Total} \\
\end{longtable}
}
\end{landscape}

\subsection{Estructura organizacional}

La estructura organizacional de Cologistic S.A.C. se presenta de forma
jerárquica y funcional, distribuida en tres niveles: estratégico,
operativo y de ejecución. En el nivel estratégico se ubican la Gerencia
General, la Subgerencia y la Gerencia de Contabilidad, responsables de
la dirección global, la administración financiera y la toma de
decisiones clave. El nivel operativo está conformado por la Jefatura de
Operaciones, el Analista de Inventario y la Jefatura de Ventas, quienes
coordinan el control de la trazabilidad y los procesos centrales de la
empresa. Finalmente, en el nivel de ejecución se encuentran los
encargados que realizan las actividades logísticas específicas como la
supervisión de despachos y picking, la operación técnica de estiba, así
como la seguridad armada y control de accesos en garita, permitiendo una
gestión segura, eficiente y alineada con los estándares de sus clientes
corporativos.

\begin{center}
  \pandocbounded{\includegraphics[width=5.90625in,height=4.46875in]{media/image7.png}}
\end{center}

\subsection{Enfoque de procesos}

En Cologistic S.A.C., la gestión organizacional se fundamenta en un
enfoque sistémico orientado a la prestación de servicios logísticos
integrales y arrendamiento de espacios a nivel corporativo (B2B).
Entendemos que la confiabilidad de nuestra cadena de suministro no
depende de acciones aisladas, sino de un flujo continuo y altamente
controlado que garantiza la trazabilidad, desde la recepción en nuestros
patios de maniobras hasta el despacho final.

Nuestra arquitectura de procesos se clasifica en tres niveles para
garantizar una operación fluida e ininterrumpida:

\begin{itemize}
\item
  \textbf{Procesos Estratégicos}: Orientan la expansión y sostenibilidad del
  negocio. Incluyen la negociación inmobiliaria para la adquisición de
  nuevas locaciones, la planificación financiera (gestión de contratos
  en dólares) y la gestión comercial B2B.
\item
  \textbf{Procesos Operativos (Misionales)}: Constituyen el \emph{core} del
  negocio e incluyen la Recepción y Despacho, la Gestión de Inventarios,
  y los Servicios Logísticos Adicionales (reempaque, repaletizado y
  \emph{picking}).
\item
  \textbf{Procesos de Soporte}: Garantizan la continuidad operativa. Abarcan la
  Gestión de Seguridad (auditorías EPP, personal armado y control de
  plagas), la Gestión de TI (implementación y mantenimiento de sistemas
  WMS/ERP) y el mantenimiento de infraestructura y maquinaria pesada.
\end{itemize}

\subsubsection{Política de Calidad y Seguridad Logística}

Cologistic S.A.C. se compromete a brindar soluciones de almacenamiento y gestión logística
confiables, seguras y adaptables a las necesidades de nuestros clientes
corporativos y estacionales. Nuestros pilares son:

\begin{enumerate}
\def\labelenumi{\arabic{enumi}.}
\item
  \textbf{Seguridad e Integridad de la Carga}: Garantizamos la protección total
  de la mercadería frente a la intemperie y riesgos externos mediante
  infraestructura adecuada y un estricto protocolo de seguridad privada
  ininterrumpida.
\item
  \textbf{Trazabilidad y Transparencia}: Nos comprometemos a mantener una
  co-gestión estrecha con nuestros clientes, asegurando la exactitud de
  los inventarios mediante auditorías físicas recurrentes y reportes
  transparentes de movimientos en tiempo real.
\item
  \textbf{Prevención y Cumplimiento Normativo}: Mantenemos un entorno de trabajo
  seguro auditando rigurosamente el uso de EPP, exigiendo el SCTR para
  todo el personal operativo, y aplicando controles sanitarios estrictos
  para el resguardo de áreas de alimentos.
\end{enumerate}

\subsection{Modelo de Negocio}

El modelo de negocio de la empresa Cologistic S.A.C., está enfocada en
ofrecer servicio de arrendamiento y logística de materias primas.

\subsubsection{Segmento de Clientes} El negocio está enfocado en el
modelo corporativo (B2B), atendiendo a grandes y medianas empresas
corporativas que manejan altos volúmenes de mercadería y requieren
espacios significativos para almacenar.

\begin{itemize}
\item
  Clientes corporativos de arrendamiento: Empresas como ISM, La Calera,
  distribuidoras de galletas (Doña Pepa) y empresas productoras de
  plásticos, las cuales buscan arrendar grandes extensiones de terreno
  (de 10,000 hasta 40,000 m²) mediante contratos a largo plazo (de 2 a 5
  años).
\item
  Clientes estacionales o "campañeros": Corporaciones como Bimbo, que
  requieren servicios de almacenamiento y gestión de operadores
  logísticos enfocados en temporadas de alta demanda (como la campaña
  navideña para panetones, de junio a diciembre). También incluye el
  almacenamiento de su área de \emph{merchandising} (material de
  marketing, folletos, góndolas).
\end{itemize}

\subsubsection{Propuesta de Valor} La empresa brinda soluciones de
almacenamiento a gran escala con flexibilidades que otros no ofrecen,
compitiendo en precio y nivel de personalización:

\begin{itemize}
\item
  Espacios seguros y resguardados: Provisión de instalaciones techadas o
  adaptadas para salvaguardar la mercadería de los clientes frente a la
  intemperie (lluvia, sol, salitre del mar).
\item
  Precios competitivos (Costos menores): Competir con gigantes del
  sector (como Ransa, Inupesa o Sabar) ofreciendo un costo menor de
  alquiler por metro cuadrado (arrendando a \$5 o \$8, en lugar de
  precios más altos).
\item
  Operador logístico integral: Para ciertos clientes estratégicos (como
  Bimbo), no solo ofrecen el espacio, sino que brindan una gestión
  logística automatizada y tercerizada que involucra: recepción,
  picking, administración de ubicaciones, control de trazabilidad,
  reempaque, repaletizado y emisión de despachos.
\end{itemize}

\subsubsection{Canales de Distribución} Dado el nivel corporativo del
servicio, la adquisición de clientes y espacios se hace por canales
físicos y directos:

\begin{itemize}
\item
  Venta y negociación B2B directa: La gerencia se encarga de ubicar
  locaciones y terrenos abandonados, negociar la gestión con los dueños
  directos y, posteriormente, buscar clientes corporativos a quienes
  ofrecer el espacio en alquiler.
\item
  Contratos físicos y notariales: El canal de cierre comercial es
  estrictamente documentado "mano a mano" y legalizado en notarías,
  donde se comprueban los títulos y la autorización de subarrendamiento.
\item
  Perfil bajo: Por medidas de seguridad ante la criminalidad, la
  captación y trato con el cliente se realiza sin intermediarios
  públicos ni letreros o logotipos en las fachadas de sus instalaciones.
\end{itemize}

\subsubsection{Relación con los Clientes} La relación está enmarcada en
la confiabilidad, formalidad y el modelo de servicio contratado:

\begin{itemize}
\item
  Relación de Independencia: Para los clientes de puro alquiler
  (terceros), la relación se basa en otorgar el espacio y dejar que el
  cliente administre su propia logística de manera independiente (como
  si fuera un departamento).
\item
  Relación Estrecha y Colaborativa (Co-gestión): Para clientes en los
  que fungen de operador logístico, la relación es intensa. Se coordinan
  inventarios semanales o quincenales cruzando información entre la
  planta del cliente y el almacén logístico, facilitando auditorías
  físicas y enviando reportes transparentes de "stock" y movimientos
  para solucionar cualquier anomalía.
\item
  Negociaciones en Dólares: Los lazos y los contratos se mantienen
  estables a través del tiempo mediante el cobro de rentas tasadas
  enteramente en dólares.
\end{itemize}

\subsubsection{Fuentes de Ingresos} El modelo monetario de la
organización proviene de dos flujos principales vinculados al uso del
espacio y la manipulación de cargas:

\begin{itemize}
\item
  Subarrendamiento de Espacios: Ingresos fuertes generados por el margen
  de ganancia entre alquilar un terreno a bajo costo (ej. \$2 el m²) y
  subarrendarlo por bloques de miles de metros cuadrados a los clientes
  finales (entre \$5 a \$8 el m²).
\item
  Tarifas de Almacenaje por Paleta: Cobros escalados según la tarifa
  acordada al mes (facturación \emph{In-Out}) basados en la cantidad de
  paletas almacenadas según el tipo de producto.
\item
  Servicios Logísticos Adicionales (VAS): Cobros adicionales por
  trabajos de manipulación física, como el "reencajado" de mercadería
  (fraccionamiento) y el "repaletizado" (trasladar cajas a paletas de
  exportación o certificadas).
\end{itemize}

\subsubsection{Recursos Clave}

\begin{itemize}
\item
  Recursos Físicos: Inmensos terrenos e infraestructura industrial en el
  Callao (almacenes y patios de maniobras de 10,000 hasta 60,000 m²).
  Maquinaria pesada como montacargas (stackers), paletas, estocas y
  barreras para el control de plagas.
\item
  Recursos Humanos: Personal certificado y de experiencia. Incluye un
  jefe de operaciones, un analista de inventarios (Ángel), de 2 a más
  montacarguistas, 3 operarios (para estiba y \emph{picking}), personal
  de limpieza y un equipo de 6 a 8 vigilantes armados para cubrir turnos
  24/7.
\item
  Recursos Tecnológicos: Actualmente limitados a computadoras y cámaras
  de seguridad. Cuentan con la urgencia de contar con un ERP o sistema
  WMS propio con lectores de código de barras para reemplazar sus
  operaciones manuales en Excel y el software desfasado "NDS".
\end{itemize}

\subsubsection{Actividades Clave}

\begin{itemize}
\item
  Negociación Inmobiliaria: Encontrar espacios subutilizados y concretar
  alquileres matrices con los terratenientes mediante procesos
  notariales para expandir la oferta de almacenes.
\item
  Recepción y Despacho: Registro exhaustivo de guías de remisión,
  control de lotes, choferes, paletas y cajas en el punto de ingreso,
  así como la planificación de las rutas de salida según la fecha de
  vencimiento.
\item
  Gestión de Inventario (Trazabilidad): Controlar movimientos internos,
  aplicar "reempaques" dividiendo paletas, ejecutar transferencias de
  códigos de una zona a otra, y consolidar reportes de \emph{kardex}
  (movimientos de ingreso y salida).
\item
  Prevención y Seguridad: Auditar el uso de EPP (Equipos de Protección
  Personal) en los almacenes, asegurar que el personal posea el SCTR
  (Seguro Complementario de Trabajo de Riesgo) y gestionar medidas
  contra roedores y palomas en áreas de alimentos.
\end{itemize}

\subsubsection{Socios Clave}

\begin{itemize}
\item
  Terratenientes y Dueños Inmobiliarios: Familias de alto poder
  económico y político en el Callao (como la familia Valle Skinner o los
  Mujica). Ellos son los verdaderos dueños de los terrenos (desde 15,000
  hasta cientos de miles de m²) a los que la empresa les alquila el área
  matriz.
\item
  Notarías Públicas: Socios legales vitales para registrar y darle
  validez jurídica a los contratos de arrendamiento ante el estado
  peruano.
\item
  Contador Externo: Profesional que recibe las órdenes y reportes de la
  empresa para emitir, digitalizar y declarar facturas ante la SUNAT.
\item
  Proveedores de Tecnología / Integradores: Desarrolladores de software
  (como el estudiante contratado originalmente para el sistema web que
  quedó obsoleto, o el actual equipo universitario) necesarios para
  construir y modernizar los sistemas operativos.
\item
  Empresas de Seguridad Privada: Proveedores externos como Prosegur y
  personal armado fundamental para la custodia de instalaciones y flujo
  de efectivo en las zonas alquiladas.
\end{itemize}
\subsection{Cadena de Valor}

La cadena de valor de Eco Logistic SAC se estructura en función de las
actividades primarias y de apoyo que permiten brindar servicios de
almacenamiento y operación logística a sus clientes. En primer lugar,
dentro de las actividades primarias, se encuentra la recepción de
mercadería, donde se realiza la validación de guías y el registro de los
productos ingresados; sin embargo, este proceso se lleva a cabo
actualmente de manera manual, lo que genera demoras y posibles errores
en el registro de información. Seguidamente, se desarrolla la actividad
de almacenamiento, que consiste en la ubicación de la mercadería dentro
de racks y el control del inventario. En esta etapa, la empresa presenta
limitaciones debido a la falta de un sistema que permita identificar con
precisión la ubicación de los productos en términos de filas, columnas y
niveles, afectando la trazabilidad.

Posteriormente, se lleva a cabo la preparación y despacho de pedidos,
que incluye actividades como el picking, embalaje y salida de
mercadería. Este proceso también presenta riesgos de error debido a la
ausencia de automatización en el control de salidas. En cuanto a la
relación con los clientes, la empresa gestiona la comunicación
principalmente a través de medios directos como correo electrónico,
WhatsApp y atención presencial, lo que evidencia una limitada
digitalización en sus canales. Asimismo, el servicio postventa se basa
en el envío de reportes en Excel y validaciones presenciales de
inventario, lo que no permite una visualización en tiempo real por parte
del cliente.

\begin{center}
  \pandocbounded{\includegraphics[width=5.90625in,height=3.51042in]{media/image3.png}}
\end{center}

\section{Situación actual del proceso a analizar}

\subsection{Matriz de Productos y Servicios}

Los cuadros presentados constituyen herramientas fundamentales dentro de
la gestión por procesos, al permitir una representación estructurada de
la relación entre la cadena de valor interna y los diferentes segmentos
de clientes de la empresa ``Productos UNIDOS''. El primer cuadro
corresponde a un listado de procesos vinculados a productos y servicios,
en el cual se identifican claramente los procesos clave (almacenamiento,
producción y comercialización), los procesos de soporte (gestión de
calidad, finanzas, recursos humanos, entre otros) y los productos
entregados al cliente, permitiendo trazar una línea directa entre la
operación y el resultado. El segundo cuadro, una matriz
producto--cliente, permite analizar la propuesta de valor desde la
perspectiva del mercado, segmentando a los clientes en mayoristas y
minoristas, precisando sus expectativas en términos de precio, calidad,
personalización y canales de atención, y estableciendo las áreas
responsables de satisfacer dichas demandas. En conjunto, ambos
instrumentos permiten fortalecer el alineamiento estratégico entre los
procesos internos y los requerimientos del entorno externo, bajo un
enfoque sistémico y orientado a la mejora continua.

\begin{center}
  \pandocbounded{\includegraphics[width=5.90522in,height=1.22222in]{media/image4.png}}
  \vspace{0.3cm}
  \pandocbounded{\includegraphics[width=5.90625in,height=2.03125in]{media/image8.png}}
\end{center}

\subsection{Matriz de Inventario de Procesos}

\begin{center}
  \pandocbounded{\includegraphics[width=5.90522in,height=4.34722in]{media/image6.png}}
\end{center}

\subsection{Matriz SIPOC del Mapa de Procesos (Nivel 0 y Nivel 1)}

\subsubsection{Nivel 0}

{\def\LTcaptype{none} % do not increment counter
\begin{longtable}[]{@{}
  >{\centering\arraybackslash}p{(\linewidth - 10\tabcolsep) * \real{0.1281}}
  >{\centering\arraybackslash}p{(\linewidth - 10\tabcolsep) * \real{0.1712}}
  >{\centering\arraybackslash}p{(\linewidth - 10\tabcolsep) * \real{0.1511}}
  >{\centering\arraybackslash}p{(\linewidth - 10\tabcolsep) * \real{0.1871}}
  >{\centering\arraybackslash}p{(\linewidth - 10\tabcolsep) * \real{0.1554}}
  >{\centering\arraybackslash}p{(\linewidth - 10\tabcolsep) * \real{0.2072}}@{}}
\toprule\noalign{}
\endhead
\bottomrule\noalign{}
\endlastfoot
& \textbf{Proveedores (del Proceso)} & \textbf{Elementos de entrada} &
\textbf{Proceso} & \textbf{Productos (Elementos de Salida)} &
\textbf{Clientes (del Proceso)} \\
\textbf{Almacenamiento} & Proveedores de materias primas, transportistas
y área de compras internas & Materias primas

orden de compra & - Recepción de materias primas

- Verificación

- Almacenamiento

- Salida hacia producción & Ingredientes pedidos para el área a
producción & Área de Producción \\
\textbf{Producción} & Área de almacenamiento & Ingredientes por

recetas & - Transformación de insumos

- Horneado

- Control de calidad

- Empaque y rotulado & Productos terminados y empacados

listos para su comercialización & Área de Ventas

Área de Logística \\
\textbf{Comercialización} & Área de producción & Productos terminados &
- Ventas

- Gestión de pedidos

- Coordinación de entregas & Órdenes de venta cerradas

Productos entregados & Clientes finales \\
\end{longtable}
}

\subsubsection{Nivel 1}

\paragraph{Almacenamiento:}

{\def\LTcaptype{none} % do not increment counter
\begin{longtable}[]{@{}
  >{\centering\arraybackslash}p{(\linewidth - 8\tabcolsep) * \real{0.1999}}
  >{\centering\arraybackslash}p{(\linewidth - 8\tabcolsep) * \real{0.2000}}
  >{\centering\arraybackslash}p{(\linewidth - 8\tabcolsep) * \real{0.2000}}
  >{\centering\arraybackslash}p{(\linewidth - 8\tabcolsep) * \real{0.2000}}
  >{\centering\arraybackslash}p{(\linewidth - 8\tabcolsep) * \real{0.2000}}@{}}
\toprule\noalign{}
\endhead
\bottomrule\noalign{}
\endlastfoot
\textbf{Proveedores (del Proceso)} & \textbf{Elementos de entrada} &
\textbf{Proceso} & \textbf{Productos (Elementos de Salida)} &
\textbf{Clientes (del Proceso)} \\
Proveedores de materias primas, transportistas y área de compras
internas & Materia prima & - Procesos de recepción prima

- Líder de Proceso: Jefe de almacenamiento & Actualización del
almacenamiento & Gerente de almacenamiento \\
Gerente de Almacenamiento & Solicitud de la salida de la materia prima &
- Revisión de ingredientes

- Separación de ingredientes

- Empaquetado por receta

- Envió de ingredientes

- Registro de ingredientes enviados & Ingredientes solicitados por
receta & Área de producción \\
\end{longtable}
}

\paragraph{Producción:}

{\def\LTcaptype{none} % do not increment counter
\begin{longtable}[]{@{}
  >{\centering\arraybackslash}p{(\linewidth - 8\tabcolsep) * \real{0.1999}}
  >{\centering\arraybackslash}p{(\linewidth - 8\tabcolsep) * \real{0.2000}}
  >{\centering\arraybackslash}p{(\linewidth - 8\tabcolsep) * \real{0.2000}}
  >{\centering\arraybackslash}p{(\linewidth - 8\tabcolsep) * \real{0.2000}}
  >{\centering\arraybackslash}p{(\linewidth - 8\tabcolsep) * \real{0.2000}}@{}}
\toprule\noalign{}
\endhead
\bottomrule\noalign{}
\endlastfoot
\textbf{Proveedores (del Proceso)} & \textbf{Elementos de entrada} &
\textbf{Proceso} & \textbf{Productos (Elementos de Salida)} &
\textbf{Clientes (del Proceso)} \\
Área de almacén & Materia prima (Ingredientes) & - Gestión de la materia
para producir masa & Masa lista y porcionada & Área de horneado \\
Encargado de Horno y cocina & Masa lista y porcionadas & - Fermentación

- Horneado

- Enfriado & Producto: Panes y postres & Área de empaquetado \\
Área de horneado & Panes y Postres & - Limpiado de producto

- Corte de productos

- Empaquetar productos & Panes y postres empaquetados & Cliente final \\
\end{longtable}
}

\paragraph{Comercialización:}

{\def\LTcaptype{none} % do not increment counter
\begin{longtable}[]{@{}
  >{\centering\arraybackslash}p{(\linewidth - 8\tabcolsep) * \real{0.1967}}
  >{\centering\arraybackslash}p{(\linewidth - 8\tabcolsep) * \real{0.1967}}
  >{\centering\arraybackslash}p{(\linewidth - 8\tabcolsep) * \real{0.1967}}
  >{\centering\arraybackslash}p{(\linewidth - 8\tabcolsep) * \real{0.2392}}
  >{\centering\arraybackslash}p{(\linewidth - 8\tabcolsep) * \real{0.1542}}@{}}
\toprule\noalign{}
\endhead
\bottomrule\noalign{}
\endlastfoot
\textbf{Proveedores (del Proceso)} & \textbf{Elementos de entrada} &
\textbf{Proceso} & \textbf{Productos (Elementos de Salida)} &
\textbf{Clientes (del Proceso)} \\
Departamento de Producción & Inventario disponible: Productos
actualizados y listos para ser distribuidos, con control preciso de
existencias.. & Confirmación y facturación: Verificación del pedido y
emisión de la factura correspondiente al cliente. & Productos
correctamente empaquetados y enviados a los clientes según las
especificaciones del pedido, asegurando que lleguen en buen estado y
dentro del plazo acordado. & Almacén \\
Departamento de Almacén & Datos de contacto y direcciones de entrega:
Información actualizada del cliente para garantizar la correcta entrega
del pedido. & Preparación y envío de productos: Embalaje y despacho de
productos según el pedido, asegurando su envío puntual. & Factura o
recibo de pago: Documento que detalla el monto total, descuentos y
productos adquiridos, confirmando el pago realizado. & Clientes
finales \\
\end{longtable}
}

\subsection{Mapa de Proceso (Nivel 0 y Nivel 1)}

\begin{center}
  \pandocbounded{\includegraphics[width=5.90625in,height=3.28125in]{media/image9.png}}
\end{center}

\end{document}
